\chapter{Topology and observers}

\section{Two slots that travel together}

This chapter does two slots: interaction topology and observer.
Together they handle the questions \emph{who interacts with whom}
and \emph{what we can actually see}. Topology and observers are not
formally connected --- a system can have a complex topology and a
trivial observer, or vice versa --- but they share a methodological
pattern: each is a structure layered on top of the substrate that the
simulator author must specify, and each is the place where most of the
implicit assumptions of a simulation live.

We will start with topology and finish with observers. The chapter
closes the four-chapter walk through the slots; Chapter~8 wires
everything together for the worked composite example.

\section{Interaction topology vs substrate topology}

The first thing to get right is the distinction between two things
that sound like the same word.

\subsection*{Substrate topology}

The \emph{substrate topology} is the intrinsic geometry of the
substrate state space $X$. For a lattice substrate, it is the lattice
$L$. For a particle substrate in $\mathbb{R}^d$, it is the manifold
$\mathbb{R}^d$. For a graph substrate, it is the graph that hosts the
function values. The substrate topology is whatever geometry the
dynamics $\Phi$ already presupposes; in the lens, the substrate
topology does \emph{not} appear in the substrate tuple $(X, \Phi,
\mu_0)$. It is implicit in $X$ and $\Phi$.

\subsection*{Interaction topology}

The \emph{interaction topology} $T$ is a different beast: it is a
(possibly weighted, possibly directed, possibly time-varying)
(hyper)graph on the \emph{entities} living on the substrate. It says
who can influence whom this tick. Where the substrate topology is
fixed by the choice of substrate, the interaction topology is a
modelling choice on top of the substrate.

In some simple cases the two coincide. A homogeneous lattice CA where
each entity is a single cell, and each cell interacts with its Moore
neighbours, has interaction topology that is identical to the
substrate lattice. This is the limiting case. For most systems of
interest, the two come apart:

\begin{itemize}[leftmargin=1.5em,itemsep=2pt]
  \item A Lenia field has lattice substrate topology, but the
    interaction topology on Lenia \emph{creatures} (entities of
    interest) is the graph of which creatures spatially overlap.
  \item A Boids swarm has manifold substrate topology
    ($\mathbb{R}^d$), but the interaction topology on boids is a
    $k$-NN graph that changes every tick.
  \item A donor-game LLM-agent population has trivial substrate
    topology, but the interaction topology is a random pairing graph
    of donors and recipients sampled per tick.
  \item A stigmergic ant simulation has lattice substrate topology
    (the floor) but the interaction topology between ants is
    \emph{indirect}: ants influence each other by leaving and reading
    pheromone trails. There is no direct ant-to-ant edge.
\end{itemize}

The lens keeps the two notions separate because they are independent
modelling axes. A change to the substrate (switching from a square to
a hexagonal lattice) is not the same kind of change as a change to
the interaction topology (switching from 4-neighbour to 8-neighbour
interaction). The first changes what the substrate \emph{is}; the
second changes what the entities \emph{do} on the substrate.

\section{Defining the interaction topology}

Formally, an \emph{interaction topology} is a measurable family
\[
  T = (T_t)_{t \in \mathbb{N}}
\]
where each $T_t$ is a hypergraph on the entity set $\mathcal{E}_t$,
together with edge weights in some semiring $R$ (typically
$\mathbb{R}_{\geq 0}$ or $\{0,1\}$). Hyperedges are subsets of entities
that interact jointly; pair hyperedges recover ordinary graphs.

Why hypergraphs rather than graphs? Because some interactions are
genuinely $k$-ary, not just pairwise. Three-way mate choice in a
parental investment model, the catalysis triplet in a RAF, the
3-clique recombination in some evolutionary algorithms. The
hypergraph generality costs nothing and accommodates those cases.

\subsection*{Common topologies}

\begin{itemize}[leftmargin=1.5em,itemsep=2pt]
  \item \textbf{Lattice neighbourhood}: Moore ($8$-cell square),
    von~Neumann ($4$-cell square), hexagonal, or higher. Standard for
    CA-like substrates.
  \item \textbf{Range-based on $\mathbb{R}^d$}: $k$-nearest-neighbour
    (each entity edges to its $k$ closest neighbours), or
    $\varepsilon$-ball (each entity edges to everything within radius
    $\varepsilon$). Standard for particle substrates.
  \item \textbf{Static social graph}: a fixed graph on agents.
    Standard for cultural-evolution models with a population
    structure that does not change.
  \item \textbf{Dynamic pairing graph}: each tick samples a fresh
    matching, e.g.\ random donor--recipient pairings. Used in
    LLM-agent donor-game systems.
  \item \textbf{Stigmergic / indirect topology}: there is no direct
    edge between entities; interactions pass through the substrate.
    An ant trail is the canonical case. Formally the topology is
    trivial; the interaction is mediated by $\Phi$ writing to and
    reading from the substrate.
\end{itemize}

\subsection*{What the weights mean}

Edge weights $w_e$ in the topology are passed to the topology-mediated
population variation kernel
\[
  V_T(\mu; x) = \sum_{e \in E(T_t)} w_e\, V(\mathrm{state}(e),
  \theta_e)
\]
that we met in Chapter~5. Weights are non-negative and typically have
units of ``mass'' --- they tell you how much each hyperedge
contributes to the resulting population. A weight of zero means the
edge is dormant for this tick. A larger weight means the edge
contributes more offspring.

\section{Multi-regime topologies: the multiplex}

Real systems often have several coupling regimes coexisting. A colony
of social insects has within-colony chemical signalling and
between-colony pheromone trails. An LLM-agent simulation has
within-cluster gossip and between-cluster broadcast. A Lenia + LLM
hybrid has spatial-overlap interactions among Lenia creatures plus
naming interactions between LLM agents and creatures.

The lens accommodates these by allowing the interaction topology to
be a \emph{layered multiplex}:
\[
  T_t = \bigsqcup_\ell T_t^{(\ell)},
\]
a disjoint union of layer-specific (hyper)graphs, each with its own
weights, plus an overall layer-weight distribution $(\alpha_\ell)$
satisfying $\sum_\ell \alpha_\ell = 1$.

The mass-conservation constraint matters. The population variation
under a multiplex topology becomes a mixture:
\[
  V_T(\mu; x) = \sum_\ell \alpha_\ell\, V_{T^{(\ell)}}^{(\ell)}(\mu;
  x),
\]
when the channels are independent. (When they are coupled through
$\diamond_\rho$, the construction generalises as in Chapter~5.) The
weights $(\alpha_\ell)$ control how much each layer contributes to the
total interaction load per tick. Without the conservation constraint,
adding a new layer would secretly increase the total amount of
variation, which would change the system's behaviour for reasons that
have nothing to do with the new coupling.

\subsection*{Multiplex example: Lenia + LLM}

Three layers:

\begin{itemize}[leftmargin=1.5em,itemsep=2pt]
  \item $T_t^{(1)}$: spatial-overlap on the Lenia lattice for Lenia
    creatures. Hyperedges = pairs of creatures whose supports
    intersect within radius $r$.
  \item $T_t^{(2)}$: uniformly random pairing of LLM judges and
    namers. Each tick, sample a matching.
  \item $T_t^{(3)}$: bipartite stigmergic layer connecting each LLM
    agent to the Lenia creatures it has named.
\end{itemize}

Weights: $\alpha_1$ controls how much Lenia-Lenia interaction
contributes, $\alpha_2$ for LLM-LLM, $\alpha_3$ for LLM-Lenia. Tuning
these is part of the modelling choice; the lens makes them visible
rather than hiding them in the simulator code.

\section{An open question: how expressive is the multiplex?}

The framework conjectures (paper OP3; see \texttt{RESULTS.md}) that
the multiplex-with-$\diamond_\rho$ formalism is universal: every
``natural'' multi-coupling-regime ALife system can be expressed as a
multiplex with countably many layers. The categorical formulation
would frame this via the operad of wiring diagrams.

OP3 is open. The book takes the multiplex as the working answer and
flags the universality question for the frontiers chapter.

\section{Observers, defined slowly}

The observer slot is the measurement apparatus. We saw the typing in
Chapter~4. Here is the longer treatment.

\subsection*{Stateless observers}

A stateless observer is a function
\[
  O : Z^w \to \mathbb{R}^k,
\]
where $Z = X \times \mathcal{M}_+(X)$ is the joint state space, $w \in
\mathbb{N} \cup \{\infty\}$ is the window length, and $\mathbb{R}^k$
holds the output. Stateless observers do not remember anything between
calls; each call evaluates a window of trajectory in isolation.

Examples:

\begin{itemize}[leftmargin=1.5em,itemsep=2pt]
  \item Mean population fitness over the last $w$ ticks.
  \item Number of distinct attractors visited in the last $w$ ticks
    (against a fixed equivalence relation).
  \item KL divergence between the current distribution and a fixed
    reference distribution.
\end{itemize}

Stateless observers are the easiest case to analyse. They are also,
in practice, the least useful: most of the interesting measurements
in artificial-life are observer-state-dependent.

\subsection*{Stateful observers}

A stateful observer carries an internal state $a_t \in A_O$ across
ticks:
\[
  O_s : Z^w \times A_O \to \mathbb{R}^k \times A_O.
\]
Each call consumes a window of trajectory and the current observer
state, produces a score vector and an updated observer state.

This is the form most actually-used observers take. The state $a_t$
might contain:

\begin{itemize}[leftmargin=1.5em,itemsep=2pt]
  \item An \textbf{archive} of previously seen behaviours (novelty
    search, Quality-Diversity).
  \item A \textbf{behaviour grid} indexing the best-performer in each
    cell (MAP-Elites).
  \item A \textbf{descriptor map} from substrate states to feature
    vectors.
  \item A \textbf{learned embedding model} that updates over time
    (some QD-with-learnt-descriptor schemes).
  \item A \textbf{reference distribution} that the observer compares
    against, possibly updated by exponential moving average.
  \item A \textbf{saturation criterion}: a rule for declaring that no
    new behaviour has been seen for some number of ticks and the run
    can stop.
\end{itemize}

The state lives entirely on the observer side. It is not part of the
substrate, not part of the population measure, not part of the
iteration rule. It is the observer's private memory of what it has
seen.

\subsection*{Read-only by default}

A subtle point: the observer does not, by default, feed back into the
iteration. The substrate dynamics, variation, viability, and topology
all run without knowing what the observer is computing. The observer
reads.

There are cases where you want feedback. Novelty search re-shapes the
variation operator using observer-state information: $V$ is biased
toward behaviours far from the current archive. The lens accommodates
this through the $\rho$-coupling apparatus of Chapter~5: the
observer's archive plays the role of an additional channel, and
$\rho_O : A_O \to \Theta$ supplies a re-parameterisation of $V$ on
first-occurrence events. The signatures of $V$ and $O_s$ remain
unchanged; the coupling is named explicitly.

\section{Three observer families}

We have promised three families. Here they are.

\subsection*{Density-distance observers}

Score a window of trajectory by a divergence between the trajectory's
empirical distribution and a reference distribution:
\[
  O[\mu_{[t-w,t]}] = D(\mu_{[t-w,t]} \,\|\, \mu_{\mathrm{ref}}).
\]
The reference $\mu_{\mathrm{ref}}$ can be a uniform distribution (for
``how spread out is the trajectory''), an archive of past behaviours
(for novelty search), or a domain-specific target distribution.

Common choices of $D$: KL divergence, Wasserstein distance, $k$-NN
distance to an archive.

This family is the one novelty search lives in. The lens makes the
common structure explicit: any divergence + any reference gives you a
density-distance observer.

\subsection*{Foundation-model embedding observers}

Embed substrate states (or trajectories) with a foundation model and
score in embedding space. ASAL embeds Lenia frames with CLIP and uses
cosine novelty against an archive of embeddings. OMNI-EPIC uses similar
embedding-space novelty for environment-population search.

These observers have an additional moving part --- the embedding
model --- which is part of the observer state. Tracking which
embedding model produced which scores matters for reproducibility.

\subsection*{Information-theoretic observers}

Score a window of trajectory by an information-theoretic quantity:
mutual information between current and future state, transfer
entropy from one part of the substrate to another, integrated
information, or --- at the macroscale --- the effective information
we used to define causal emergence in Chapter~3.

These observers are the most demanding computationally: they require
estimating distributions over the substrate state space. They are
also the most rigorous: they have a clear meaning that does not
depend on a learned model or an arbitrary descriptor space.

\section{Open conjecture: are the three families one?}

Conjecture~C3 (see \texttt{RESULTS.md}): \emph{the three observer
families are all instances of a single divergence-functional schema}
\[
  O[\mu] = D(\mu \,\|\, \mu_{\mathrm{ref}})
\]
\emph{for a context-dependent reference $\mu_{\mathrm{ref}}$.}

Density-distance observers fit by construction. Foundation-model
observers fit if you let $\mu_{\mathrm{ref}}$ be the archive's
embedding distribution. Information-theoretic observers fit if you
let $\mu_{\mathrm{ref}}$ be the appropriate factored distribution
(e.g.\ the product of marginals for mutual information).

C3 is plausible by inspection but unproven in generality. The lens
takes the three families as separate; if C3 is proved, the families
unify.

\section{Window length, archive size, and saturation}

Three observer-side parameters that recur in every implementation.

\paragraph{Window length $w$.} How many ticks of trajectory the
observer consumes per evaluation. Short windows respond quickly to
transient behaviour; long windows are robust to noise. The lens
treats $w$ as part of the observer specification --- a system with
the same observer family but a different window is a different
system. The system description should record $w$.

\paragraph{Archive size $A$.} For density-distance and FM-embedding
observers, the archive size is the cap on how many past behaviours
are stored. Larger archives are more expressive but more expensive
per evaluation (every new behaviour is compared against the archive).
Pruning strategies (random removal, density-based removal,
behavioural-distance-based removal) matter.

\paragraph{Saturation.} A trajectory is said to have saturated if no
new behaviour has been added to the archive for some number of ticks.
Saturation is the practical sign that a system has run out of
novelty; it does not necessarily mean the system has ``stopped'' in
any substrate-dynamics sense.

\section{Worked example: the Lenia + LLM observer}

For the Lenia + LLM hybrid we have been building up, a natural
observer is the pair
\[
  O = (O_\Omega,\, O_{\mathrm{CLIP}})
\]
with a shared window $w = W$:

\begin{itemize}[leftmargin=1.5em,itemsep=2pt]
  \item $O_\Omega$ is the residence-time $\Omega$ metric on the time
    series of distinct Lenia attractors observed in the window.
    Information-theoretic family.
  \item $O_{\mathrm{CLIP}}$ is the mean cosine novelty of CLIP
    embeddings of (frame, name) pairs against an archive of size
    $A$. FM-embedding family.
\end{itemize}

The observer reports a two-dimensional score vector each tick. The
observer state contains the attractor history (for $O_\Omega$) and the
CLIP archive (for $O_{\mathrm{CLIP}}$). Neither feeds back into the
iteration by default; the system runs forward regardless of what the
observer sees.

\section{What we have now}

After four chapters, the lens is fully built. To recap:

\begin{itemize}[leftmargin=1.5em,itemsep=2pt]
  \item The substrate $\mathbf{X} = (X, \Phi, \mu_0)$ holds the
    material and its native dynamics. Composes via $\boxtimes_\kappa$.
  \item Variation $V : X \times \Theta \to \mathcal{P}(X)$ proposes
    candidates. Composes via $\diamond_\rho$.
  \item Viability $F : \mathcal{M}_+(X) \to \mathcal{M}_+(X)$ filters.
    Four formalisms; closure-operator unification is open.
  \item Topology $T = (T_t)_{t \geq 0}$ describes who-affects-whom.
    Multi-regime via multiplex; universality of multiplex is open.
  \item Observer $O_s$ measures. Three families; unification is open.
\end{itemize}

The next chapter brings them together for the worked composite
example and the quantitative measurement of emergence and novelty.

\section*{To think about}

\begin{enumerate}
  \item Pick a system where the substrate topology and the
    interaction topology come apart most dramatically. What does that
    decoupling enable?
  \item The multiplex formalism requires the layer weights
    $(\alpha_\ell)$ to sum to one. What goes wrong if they do not?
  \item Novelty search uses observer feedback to bias variation. The
    lens handles this with $\rho$-coupling. Sketch how an MCC system
    (Chapter~6) would similarly couple observer state to variation,
    if at all.
  \item C3 conjectures that the three observer families unify under
    divergence-from-reference. What would a counterexample look like?
  \item Saturation is the practical signal that a system has run out
    of novelty. Is saturation always a property of the observer
    rather than the substrate?
\end{enumerate}
